\section{Comparación de familias y problemas pendientes}
\label{sec:cobertura-cuatro-familias}

La comparación comprende dos características independientes: la estructura
del campo y el orden de la derivada. Los sistemas tipo Chua conservan la
arquitectura del circuito con una no linealidad escalar; los sistemas no Chua
incluyen tanto otras realizaciones de Lur'e como campos con varias
interacciones no lineales. La representación de Lur'e, por sí sola, no
identifica un sistema como Chua. El orden entero corresponde a \(q=1\); el
orden fraccionario conmensurable, a una derivada de Caputo común de orden
\(0<q<1\).

En cada familia se distinguen cuatro problemas: resolver las ecuaciones
algebraicas, construir una semilla, determinar su destino dinámico y
caracterizar la cuenca del conjunto alcanzado. La existencia de equilibrios o
de puntos perpetuos resueltos no implica la existencia de un atractor
localizado. Asimismo, un resultado negativo, como la ausencia de puntos
perpetuos o la convergencia hacia un equilibrio, delimita la aplicabilidad
del método y forma parte de la comparación.

\subsection{Resultados en las cuatro familias}

\begin{longtable}{@{}L{0.16\linewidth}L{0.25\linewidth}L{0.25\linewidth}L{0.26\linewidth}@{}}
\caption{Cobertura algebraica y dinámica por familia.}
\label{tab:cobertura-cuatro-familias}\\
\toprule
Familia & Sistemas representativos & Desarrollo disponible & Problema pendiente\\
\midrule
\endfirsthead
\toprule
Familia & Sistemas representativos & Desarrollo disponible & Problema pendiente\\
\midrule
\endhead
Chua entero
& Chua PWL canónico; arctangente como sistema auxiliar de búsqueda.
& Equilibrios, espectros, transferencia, función descriptiva y semilla modal
del PWL; continuación, diagnósticos dinámicos y comparación independiente.
& Extender la reproducción a las ramas coexistentes no evaluadas con el
mismo detalle. La exclusión de vecindades continuas sigue sin demostrarse.\\
No Chua entero
& Kalman--Fitts, MAVPD y PLL.
& Álgebra y localización numérica; ciclos de Kalman--Fitts, MAVPD y PLL;
candidato caótico de MAVPD. Sondeos locales y comparación entre integradores.
& Certificar atracción y separación de cuencas; completar el análisis de
frontera. Los puntos perpetuos no sustituyen las otras rutas.\\
Chua fraccionario
& PWL de Danca, arctangente de Wu y c590.
& Álgebra, estabilidad de Matignon y semillas; continuación causal PWL;
candidato c590; trayectorias desde semillas FPP--A de Wu.
& Refinamiento de paso del extremo PWL sesgado; destino asintótico de las
semillas de Wu; convergencia conjunta de memoria, horizonte y clasificación.\\
No Chua fraccionario
& Lorenz generalizado, Rabinovich--Fabrikant y Muñoz--Pacheco.
& Realizaciones, equilibrios y semillas FPP--A; absorción de Lorenz;
trayectorias preliminares de Muñoz--Pacheco.
& Integrar y clasificar Lorenz y RF desde las semillas calculadas; completar
convergencia, atracción y clasificación del caso Muñoz--Pacheco.\\
\bottomrule
\end{longtable}

Las cuatro familias cuentan con resultados algebraicos. La cobertura dinámica
es desigual: los sistemas enteros poseen varios casos de localización y los
Chua fraccionarios cuentan con candidatos y controles extensos. En los
sistemas fraccionarios no Chua, el resultado dinámico de Muñoz--Pacheco es
preliminar; las pruebas de absorción de Lorenz y las raíces de
Rabinovich--Fabrikant aún no se acompañan de una caracterización local de sus
atractores. Los cálculos de exponentes descritos en la
\cref{sec:nonchua-dynamic-benchmarks} son comparaciones de métodos y no
completan esa localización.

\subsection{Procedimientos pendientes por sistema}

\begin{longtable}{@{}L{0.24\linewidth}L{0.32\linewidth}L{0.36\linewidth}@{}}
\caption{Desarrollos necesarios para completar la comparación.}
\label{tab:procedimientos-pendientes}\\
\toprule
Sistema & Resultado disponible & Desarrollo pendiente\\
\midrule
\endfirsthead
\toprule
Sistema & Resultado disponible & Desarrollo pendiente\\
\midrule
\endhead
Chua PWL entero
& Semilla armónica y candidato caótico canónico; ausencia de PP regionales.
& Reproducir las otras ramas caótica y periódica de coexistencia con
condiciones iniciales, paso y horizonte explícitos; comparar sus cuencas.\\
Chua PWL, \(q=0.9998\)
& Dos semillas centradas; extremo sesgado regular; 36 perturbaciones del dato
inicial a paso fijo.
& Evaluar \(h,h/2,h/4\) con memoria completa y horizonte común; distinguir
sensibilidad al dato inicial de convergencia del integrador.\\
Chua arctangente de Wu, \(q=0.99\)
& Equilibrios, cierre armónico, dos FPP--A y trayectorias de dos integradores.
& Resolver el destino de los transitorios mediante refinamiento y ampliación
del horizonte. La recurrencia ADM publicada y el PVI Caputo causal requieren
comparaciones separadas.\\
Chua arctangente c590, \(q=0.9999\)
& Candidato aperiódico y sondeos locales; un par FPP--A alejado de la
trayectoria candidata.
& Evaluar ese par con controles de escape y perturbaciones; demostrar la
separación de vecindades si se pretende una clasificación rigurosa.\\
Kalman--Fitts, \(q=1\)
& Ciclo estable numérico, ruta de conmutación y obstrucción del cierre directo;
ausencia de PP.
& Establecer un dominio de retorno y una cota de contracción, o un argumento
equivalente de atracción; calcular una vecindad explícita del origen contenida
en su cuenca local de estabilidad.\\
MAVPD, \(q=1\)
& PP exactos y ciclos para \(\xi=3.1\); control con contacto para
\(\xi=3.5\); candidato caótico para \(\xi=2.85\).
& Comparar PP, función descriptiva y búsqueda directa bajo los mismos
parámetros; caracterizar las tres cuencas detectadas al aproximar la frontera.\\
PLL, \(q=1\)
& Órbita corredora en el cilindro y dos PP por vuelta; un PP llega al
equilibrio y el otro permanece transitorio en el experimento geométrico.
& Resolver el segundo destino, certificar el retorno cilíndrico mediante
encierros intervalares y caracterizar la componente de frontera detectada.\\
Lorenz generalizado, \(q=0.995\)
& Tres equilibrios, dos FPP--A, balance MIMO, existencia global, absorción y
precompacidad de perfiles.
& Integrar las dos semillas y controles próximos; resolver el balance MIMO si
se usa esa ruta; determinar recurrencia, atracción y cuencas.\\
Rabinovich--Fabrikant, \(q=0.998\), \((a,b)=(0.1,0.2876)\)
& Cinco equilibrios, ocho FPP--A y formulación de balance MIMO.
& Integrar los cuatro pares de semillas; controlar existencia y acotación,
comparar métodos y horizontes, y estudiar vecindades de los cinco equilibrios.\\
Muñoz--Pacheco, \(q=0.97\)
& Ausencia de equilibrios, cuatro FPP--A en regiones suaves y trayectorias
hasta \(t=35\); una recta de candidatos direccionales en la frontera genera
soluciones no acotadas.
& Obtener convergencia de paso y horizonte, acotación y atracción; contrastar
las semillas con los datos publicados. El evento FPP--H no está resuelto.\\
Muñoz--Pacheco, \(q=0.996\)
& Parámetros e iniciales publicados para coexistencia.
& Reproducir las dos respuestas con memoria completa y comprobar la
coexistencia bajo el mismo método, tolerancias y horizonte.\\
Jerk cuadrático de Liu
& Campo, realización, equilibrios y ausencia de PP no estacionarios.
& Localizar mediante balance, continuación o búsqueda directa; estudiar
destino, robustez y cuenca.\\
Liu extendido
& Realización MIMO exacta y semirrecta de equilibrios
\(x=y=z=0,\ w<0\), con \(\operatorname{sign}(0)=0\).
& Fijar el orden y la convención de solución en la discontinuidad; estudiar
estabilidad de la variedad de equilibrios y resolver PP por regiones antes
de construir e integrar semillas.\\
\bottomrule
\end{longtable}

Para Lorenz estándar, RF con \(b=0.98\), jerk exponencial, sistema financiero
y sistema de cuatro alas existen desarrollos algebraicos y comparaciones de
exponentes. Los parámetros de estos ensayos son distintos de los casos de
atractores ocultos de Lorenz generalizado y RF con \(b=0.2876\).
Su incorporación como casos de localización requiere estudiar semillas y
cuencas de manera específica. Las discrepancias entre exponentes deben
resolverse antes de utilizarlos como referencia cuantitativa.

Los candidatos adicionales de la \cref{sec:candidate-queue} requieren datos
previos al cálculo dinámico. En Genesio--Tesi faltan los parámetros y las
condiciones iniciales del caso seleccionado; en el oscilador memristivo, la
constante de la variedad invariante y sus parámetros; en el jerk PWA, el
campo regional y las superficies de conmutación. El ejemplo (162) de
Leonov--Kuznetsov requiere precisar su característica PWA, y el jerk con seno
hiperbólico exige resolver la ambigüedad entre esa función y su inversa en la
fuente. Sus representaciones estructurales constituyen puntos de partida
algebraicos, sin resultados locales de clasificación de atractores.

\subsection{Cobertura del procedimiento de puntos perpetuos}

Para un campo \(f\in C^1\), la resolución algebraica consiste en determinar
\[
 \mathcal P(f)=\{p:J_f(p)f(p)=0,\ f(p)\ne0\}.
\]
La condición \(\det J_f(p)=0\) reduce la búsqueda, pero no sustituye las
ecuaciones vectoriales. En campos definidos por regiones se resuelve cada
sistema interior, se comprueba la pertenencia a su región y se estudian por
separado las superficies donde el jacobiano no existe. Una enumeración
completa exige justificar la ausencia de otras raíces reales, mediante
eliminación algebraica, aislamiento de raíces o una demostración directa.

La \cref{sec:pp-case-audit} proporciona raíces para Chua arctangente, MAVPD,
PLL, Lorenz generalizado, RF con \(b=0.2876\) y Muñoz--Pacheco, junto con
controles de ausencia para Chua PWL y Kalman--Fitts. Los dos sistemas jerk
también admiten una demostración de ausencia. Permanecen pendientes los
desarrollos específicos de PP para Lorenz estándar, RF con \(b=0.98\),
sistema financiero, sistema de cuatro alas y Liu extendido. La ausencia de PP
en un sistema no impide localizar atractores mediante otros métodos.

En Muñoz--Pacheco, las cuatro raíces suaves no incluyen la recta
\(\{(0,y,0):y\in\mathbb R\}\), situada en la superficie de conmutación.
Allí no existe un jacobiano clásico, aunque ambos jacobianos laterales
anulan el campo. La solución exacta permanece en esa recta y satisface
\(y(t)=y_0+(t-t_0)^q/\Gamma(q+1)\); su crecimiento excluye la atracción
por un conjunto acotado. Esta rama pertenece al criterio direccional por
regiones y debe distinguirse de los puntos perpetuos clásicos.

Si se conserva el mismo campo y sólo cambia \(q\), las raíces de
\(J_ff=0\) permanecen iguales. En orden entero son puntos de aceleración
nula; para Caputo conmensurable definen semillas FPP--A mediante la expansión
de arranque. El orden modifica la trayectoria y su estabilidad, por lo cual
la evaluación dinámica debe repetirse. El evento FPP--H depende de la
historia y requiere un cálculo adicional del operador secuencial y de su
convergencia.

\subsection{Función del análisis topológico}

La topología interviene en problemas distintos. El grado de Brouwer permite
demostrar la existencia de un equilibrio; aplicado a \(I-P\), donde \(P\)
es un retorno bien definido, puede establecer la existencia de una órbita
periódica. Las variedades invariantes y las separatrices orientan la búsqueda
de fronteras de cuenca. Un bloque aislante delimita dinámica invariante y
permite aplicar el índice de Conley. Cada herramienta requiere hipótesis
propias de continuidad, compacidad, transversalidad o aislamiento, desarrolladas
en las \cref{sec:geo-degree,sec:geo-poincare-floquet,sec:metodos-topologicos-existencia,sec:localizacion-geometrico-topologica}.
El ejemplo de la \cref{subsec:guide-hidden-cycle-topology} combina un anillo
atrapante, un retorno y una separatriz para demostrar la existencia de un
ciclo oculto y calcular su cuenca exactamente.

La clasificación de ocultedad añade una condición sobre la cuenca. Una vez
establecidos el atractor \(\mathcal A\) y su espacio de fases \(X\), deben
existir vecindades \(U_e\) de todos los equilibrios tales que
\[
 U_e\cap\mathcal B_X(\mathcal A)=\varnothing.
\]
El grado y el índice local de un equilibrio no implican esta exclusión. En
el caso entero, una región positivamente invariante contenida en otra cuenca
puede proporcionar el argumento. Para Caputo, la condición debe formularse
en un espacio de perfiles admisibles o respecto de una clase de datos
iniciales definida; la proyección puntual no conserva toda la información.

En Kalman--Fitts, la estabilidad asintótica del origen ya garantiza la
existencia de una vecindad atraída por él. Si se demuestra un atractor
compacto separado del origen, esa vecindad queda excluida de su cuenca. El
cálculo de un radio explícito y la certificación del ciclo son pasos
adicionales.

Los cálculos de la \cref{sec:resultados-piloto-gt} contienen aproximaciones
de frontera y destinos de semillas geométricas. El intervalo de PLL y el
tercer destino de MAVPD orientan el muestreo posterior, pero aún no constituyen
una caracterización topológica de la frontera. Las superficies críticas y
los PP restringen regiones de búsqueda. La certificación puede apoyarse en
un retorno, un conjunto aislante u otro argumento que demuestre atracción
y separación de cuencas.
